\scanpage{96}
\mspage{\handwritten{$2'$}}
and for $\ell=1$
\[
\begin{aligned}
s_1(v*_1u)&=s_1(u)\\
t_1(v*_1u)&=t_1(v)
\end{aligned}
\]
(NB the operation $v*_1u$ is the usual composition $v\circ u$).

One may be tempted to think that the preceding data exhaust the structure of $\infty$-groupoid, and that they will have to be supplemented only by a handful of suitable \uline{axioms}, one being \uline{associativity} for the operation $\overset{i}{*}_\ell$, which can be expressed essentially by saying that this composition operation turns $F_i$ into the set of arrows of a category having $F_{i-\ell}$ as set of objects (with\deleted{\unsure{s}} the source and target maps $s^i_\ell$ and $t^i_\ell$, and with identity map $k_\ell^{i-\ell}:F_{i-\ell}\longrightarrow F_i$, the composition of the identity maps $F_{i-\ell}\longrightarrow F_{i-\ell+1}\longrightarrow\ldots F_{i-1}\longrightarrow F_i$), and another being the Godement relation
\[
(v' *_{\handwritten{\alpha}}v)*_{\handwritten{\nu}}(u'*_{\handwritten{\alpha}}u)
=(v'*_{\handwritten{\nu}}u')*_{\handwritten{\alpha}}(v*_{\handwritten{\nu}}u)
\]
(with the assumptions \handwritten{$1\leq\alpha<\nu$,} \deleted{\unsure{and}} \handwritten{$u,u',v,v'\in F_i$ and}
\[
\handwritten{\left\{\begin{aligned}t_\alpha(u)&=s_\alpha(u')\\t_\alpha(v)&=s_\alpha(v')\end{aligned}\right. \qquad
t_\nu(u)=s_\nu(v)=s_\nu(v')=t_\nu(u')}
\]
\deleted{\unsure{imply}}\handwritten{\unsure{express}}ing that both members are defined), plus the two relations concerning the inversion of $i$-objects $(i\geq1)$ $u\mapsto\check{u}$,
\[
\begin{aligned}
u*_1\check{u}&=\deleted{\unsure{k_1}}\ \mathrm{id}_{t_1(u)}\ ,\qquad \check{u}*_1u=\mathrm{id}_{s_1(u)}\\
{(v*_\ell u)}^{\scriptscriptstyle\vee}&=\check{\unsure{v|u}}*_\ell\check{u}\qquad(\ell\geq2)\quad.
\end{aligned}
\]
It just occurs to me, by the way, that the \inserted{previous} description of basic (or ``primary'') data for an $\infty$-groupoid is already uncomplete in some rather obvious respect, namely that the symmetry-operation \inserted{\handwritten{$\mathrm{inv}_i:$}} $u\mapsto\check{u}$ on $F_i$ must be complemented by $i-1$ similar involutions on $F_i$, which corresponds algebraically to the intuition that when we have an \inserted{\handwritten{$(i+1)$-}}arrow $\handwritten{\lambda}$ say between two \inserted{\handwritten{$i$-}}arrows $u$ and $v$, then we must be able to deduce from it another arrow from $\check{u}$ to $\check{v}$ (namely $u\mapsto\check{u}$ has a ``functorial \unsure{caracter|character}'' \deleted{\unsure{in u}} for variable $u$)\unsure{?|;|.} This seems a rather anodine modification of the previous set-up, and is irrelevant for the main point I want to make here, namely: that for the notion of $\infty$-groupoids we are after, all the equalities just envisioned in this paragraph (and those I guess which will ensue\deleted{s} naturally by the necessary extension of the basic involution on $F_i$) should be replaced by ``homotopies'' namely by $(i+1)$-arrows between the\deleted{\unsure{r}} two members. These arrows should be viewed, I believe, as being part of the data, \deleted{the} \deleted{The diffi} they appear here as a kind of ``secondary'' structure. The difficulty which appears now is to work out the natural coherence properties
% archive sheet number: 95
